% Compilar com: latexmk -pdf Relatorio_Atividade_2_ABNT.tex
\documentclass[12pt,a4paper]{article}

\usepackage[utf8]{inputenc}
\usepackage[T1]{fontenc}
\IfFileExists{brazil.ldf}{\usepackage[brazil]{babel}}{\usepackage[english]{babel}\frenchspacing}
\usepackage{mathptmx}
\usepackage{microtype}
\usepackage{geometry}
\usepackage{setspace}
\usepackage{indentfirst}
\usepackage{titlesec}
\usepackage{booktabs}
\usepackage{tabularx}
\usepackage{array}
\usepackage{float}
\usepackage{needspace}
\usepackage{listings}
\usepackage{xcolor}
\usepackage[hidelinks]{hyperref}

\AtBeginDocument{%
  \renewcommand{\tablename}{Tabela}%
  \renewcommand{\refname}{Referências}%
}

\geometry{top=3cm,left=3cm,right=2cm,bottom=2cm}
\setlength{\parindent}{1.25cm}
\setlength{\parskip}{0pt}
\onehalfspacing
\emergencystretch=2em
\hyphenpenalty=8000

\titleformat{\section}
  {\normalfont\bfseries\MakeUppercase}{\thesection}{0.7em}{}
\titlespacing*{\section}{0pt}{1.2\baselineskip}{0.5\baselineskip}

\newcolumntype{Y}{>{\raggedright\arraybackslash}X}
\renewcommand{\arraystretch}{1.12}

\definecolor{codigoFundo}{RGB}{246,247,249}
\definecolor{codigoBorda}{RGB}{205,210,217}
\lstset{
  basicstyle=\ttfamily\small,
  backgroundcolor=\color{codigoFundo},
  frame=single,
  rulecolor=\color{codigoBorda},
  columns=fullflexible,
  keepspaces=true,
  showstringspaces=false,
  breaklines=true
}

\hypersetup{
  pdftitle={Analisador Léxico para Mini Pascal},
  pdfauthor={João Paulo Sandes Brito},
  pdfsubject={Atividade 2 da disciplina de Compiladores},
  pdfkeywords={analisador léxico, Java, Mini Pascal}
}

\begin{document}

\begin{center}
  {\bfseries\MakeUppercase{Analisador léxico para a linguagem Mini\_Pascal}}\par
  \vspace{0.8\baselineskip}
  João Paulo Sandes Brito\par
  \small Universidade Estadual do Sudoeste da Bahia -- UESB\par
\end{center}

\vspace{0.5\baselineskip}

\renewcommand{\abstractname}{Resumo}
\begin{abstract}
\noindent
O trabalho apresenta um analisador léxico para Mini\_Pascal. O programa foi escrito em Java 17, possui interface em Swing e gera pares no formato lexema--token. Foram usados seis arquivos de teste, incluindo um conjunto com erros léxicos.

\vspace{0.3\baselineskip}
\noindent\textbf{Palavras-chave}: análise léxica; Java; Mini\_Pascal.
\end{abstract}

\begin{center}
  \textbf{LEXICAL ANALYZER FOR THE MINI\_PASCAL LANGUAGE}
\end{center}

\renewcommand{\abstractname}{Abstract}
\begin{abstract}
\noindent
This paper presents a lexical analyzer for Mini\_Pascal. The program was written in Java 17, provides a Swing interface and generates lexeme--token pairs. Six input files were used for testing, including lexical errors.

\vspace{0.3\baselineskip}
\noindent\textbf{Keywords}: lexical analysis; Java; Mini\_Pascal.
\end{abstract}

\section{Objetivo}

O analisador lê um arquivo-fonte e separa seu conteúdo em tokens. Espaços, quebras de linha e comentários são ignorados. O resultado pode ser visto na interface ou salvo em arquivo texto.

A implementação usa transições determinísticas e mantém a linha e a coluna de cada lexema. Essa organização corresponde ao funcionamento usual da etapa léxica de um compilador \cite{aho2007}.

\section{Classes léxicas}

\begin{table}[H]
\centering
\caption{Classes reconhecidas}
\small
\begin{tabularx}{\textwidth}{p{3.1cm}Y Y}
\toprule
\textbf{Classe} & \textbf{Exemplos} & \textbf{Token} \\
\midrule
Palavra reservada & \texttt{program}, \texttt{begin}, \texttt{integer} & Palavra reservada \\
Identificador & \texttt{x}, \texttt{var2}, \texttt{media\_geral} & Identificador \\
Número & \texttt{13}, \texttt{1.33}, \texttt{24.40e-04} & Inteiro ou real \\
Constante & \texttt{'A'}, \texttt{"texto"} & Caractere ou string \\
Operador & \texttt{+}, \texttt{mod}, \texttt{>=}, \texttt{and} & Aritmético, relacional ou lógico \\
Símbolo & \texttt{( ) , ; :}, \texttt{:=}, \texttt{.} & Especial, atribuição ou fim \\
\bottomrule
\end{tabularx}
\end{table}

Palavras reservadas não diferenciam maiúsculas de minúsculas, mas o lexema original é mantido. Os termos \texttt{read}, \texttt{write} e \texttt{writeln} seguem o exemplo da atividade e são classificados como identificadores. O símbolo \texttt{=} é tratado como operador relacional.

\section{Implementação}

\begin{table}[H]
\centering
\caption{Principais classes do programa}
\small
\begin{tabularx}{\textwidth}{p{4cm}Y}
\toprule
\textbf{Classe} & \textbf{Função} \\
\midrule
\texttt{Lexer} & Reconhecer lexemas e produzir tokens. \\
\texttt{Token/TokenType} & Guardar lexema, tipo, linha, coluna e mensagem. \\
\texttt{AnalysisFileService} & Ler a entrada e gravar a saída. \\
\texttt{MainFrame} & Exibir a interface gráfica. \\
\texttt{Main} & Iniciar o programa. \\
\bottomrule
\end{tabularx}
\end{table}

O método \texttt{nextToken()} lê um token por chamada. O método \texttt{scanAll()} repete a operação até o fim do arquivo. A decisão inicial depende do primeiro caractere do lexema:

\begin{lstlisting}
letra ou '_' -> identificador ou palavra reservada
digito       -> numero inteiro ou real
aspas        -> constante
operador     -> forma simples ou composta
outro        -> erro lexico
\end{lstlisting}

\section{Erros léxicos}

Os erros registram lexema, linha, coluna e mensagem. Depois do registro, a leitura continua no próximo ponto válido. São detectados símbolos desconhecidos, expoentes incompletos, números unidos a identificadores, comentários ou strings sem fechamento e constantes de caractere inválidas.

\Needspace{16\baselineskip}
\section{Testes}

\begin{table}[H]
\centering
\caption{Casos executados}
\small
\begin{tabularx}{\textwidth}{p{2cm}Y c c}
\toprule
\textbf{Caso} & \textbf{Conteúdo} & \textbf{Tokens} & \textbf{Erros} \\
\midrule
Teste 01 & Programa básico & 50 & 0 \\
Teste 02 & Números e constantes & 48 & 0 \\
Teste 03 & Operadores & 74 & 0 \\
Teste 04 & Comentários & 15 & 0 \\
Teste 05 & Entradas inválidas & 30 & 5 \\
Teste 06 & Maiúsculas, sublinhado e Unicode & 32 & 0 \\
\bottomrule
\end{tabularx}
\end{table}

Os seis arquivos produziram 249 tokens. As saídas foram comparadas com os arquivos esperados. O teste automatizado também executou vinte verificações.

\section{Execução}

No Windows, o programa pode ser aberto por \texttt{executar.bat}. O JAR também aceita entrada e saída pela linha de comando:

\begin{lstlisting}
java -jar dist\MiniPascal-Lexico.jar entrada.txt saida.txt
\end{lstlisting}

É necessário ter Java 17 ou superior. Não é preciso usar uma IDE.

\section{Considerações finais}

O programa reconhece as classes previstas, informa erros sem interromper toda a leitura e gera a listagem lexema--token. Os arquivos de teste e o JAR acompanham o código-fonte.

\begin{thebibliography}{9}
\setlength{\itemsep}{0.4\baselineskip}

\bibitem{aho2007}
AHO, Alfred V.; LAM, Monica S.; SETHI, Ravi; ULLMAN, Jeffrey D. \textbf{Compilers: principles, techniques, and tools}. 2. ed. Boston: Pearson Addison-Wesley, 2007.

\end{thebibliography}

\end{document}
